% SPDX-License-Identifier: CC-BY-4.0
% =====================================================================
% Section 3 of D. Joyner, "A hitch-hiker's guide to the metaplectic trace
% formula for SL(2)", revised for the 2-fold cover (n = 2), November 2026.
%
% Stage 2 of ../revision_log.md. Source: ../original-1995/cctrace2.tex
% (= ../book-1997/book7b.tex). For n = 2 the cover of a torus is abelian, so
% the Clifford-Mackey part of the 1995 section (Lemmas 3.3, 3.5, Remarks
% 3.4, 3.6, 3.7, 3.10) is not needed; the general-n situation is summarized
% in Remark 2026.6. The map from 1995 numbers to the numbers here is in the
% "Stage 2 notes" of ../revision_log.md. Explicit formulas are checked by
% checks/sec3_n2_check.py.
% Revised after the second referee report (ChatGPT, November 2026): test
% functions are defined on the topological cover, and the proof of
% Lemma 3.1(c) no longer claims that a continuous character has open kernel.
% Comments beginning "% [Claude]" are editorial notes by Claude.
% =====================================================================

\section{The rank 0 case}\label{sec:rank0}

Before the trace formula for $\overline{G(\A)}$ we treat the covers of
the tori of $G$, where it is elementary. For $n=2$ these covers are
abelian (Lemma~\ref{lem:global}(e)), but in general not split, so a
genuine automorphic character cannot be written as an ordinary one
times a sign (Example~\ref{ex:order4}). The trace formula is then
Poisson summation twisted by one fixed genuine automorphic character
(Theorem~\ref{thm:tf-rank0}, Lemma~\ref{lem:transfer}), and it gives a
bijection between genuine automorphic characters and Hecke characters
(Theorem~\ref{thm:corr}). For the split torus, Weil's character
$\chi_\psi$ (Lemma~\ref{lem:weil}) makes the bijection canonical once
$\psi$ is fixed. We also define the distributions $I^A_{\mathcal O}$ and
$I^A_{\bar\chi}$ used in Section~6. We work with an arbitrary maximal
torus because Section~7 applies the same formula to anisotropic tori.

\subsection{The cover of a torus}\label{sec:torus-cover}

Let $T$ be a maximal torus of $G$ defined over $\Q$. Then
$T=\Res_{k/\Q}T_k$ for a maximal torus $T_k$ of $SL(2)$ over $k$, which
is either split, and then $GL(2,k)$-conjugate to $A$, or of the form
$T_\theta$ of Section~\ref{sec:local} with $E=k(\sqrt\theta)$ a quadratic
extension of $k$. Put
\[
T(\A)^1=\{t\in T(\A)\mid |\chi(t)|=1\ \text{for all}\ \chi\in X(T)_\Q\}.
\]
For $T=A$ this is $A(\A)^1=\{a(t)\mid |t|_{\A_k}=1\}$, the kernel of
$H$ of Lemma~\ref{lem:NAK-global}; for $T_k=T_\theta$ the torus $T$ is
anisotropic over $\Q$ and $T(\A)^1=T(\A)$. In both cases
$T(\Q)\subset T(\A)^1$ and $T(\Q)\backslash T(\A)^1$ is
compact~\cite{PR}.

Let $\overline{T(\A)}$ and $\overline{T(\A)^1}$ be the preimages in
$\overline{G(\A)}$, and $\overline{T(k_v)}$ the preimage of $T_k(k_v)$
in $\overline{SL(2,k_v)}$. By Lemmas~\ref{lem:local}(g)
and~\ref{lem:global}(e) these groups are abelian. On $A$ the cocycle is
\begin{equation}\label{eq:cocycle-A}
\sigma_\A(a(t_1),a(t_2))=\prod_v(t_{1,v},t_{2,v})_v,\qquad
t_1,t_2\in\A_k^\times,
\end{equation}
since $\kappa_v\equiv1$ on $A$. Exactly as for $G$ in
Section~\ref{sec:global}, $\overline{T(\A)}$ is the \emph{quotient} of the
restricted product of the $\overline{T(k_v)}$, taken with respect to the
subgroups $T(\OO_v)^*=\{(t,\kappa_v(t))\mid t\in T_k(k_v)\cap K_v\}$, by
$\{(\zeta_v)\mid\prod_v\zeta_v=1\}$; it is not a subgroup of that
restricted product. Each $\overline{T(k_v)}$ maps injectively into
$\overline{T(\A)}$. These are closed subgroups of the covers of
Section~\ref{sec:background}, and they carry the induced topologies; for
$\overline{T(\A)}$ this is again the quotient topology of the restricted
product. We regard $T(\Q)$ as a subgroup of $\overline{T(\A)^1}$
by Lemma~\ref{lem:global}(a), $\gamma\mapsto(\gamma,\kappa_\A(\gamma))$;
for $T=A$ this is $a(t)\mapsto(a(t),1)$, $t\in k^\times$.

\paragraph{Conventions.}
As in Section~\ref{sec:functions}, a representation $\pi$ of one of these
covers is \emph{genuine} if $\pi(1,\zeta)=\zeta$, and a function $f$ is
\emph{anti-genuine} if $f(x\zeta)=\zeta^{-1}f(x)$ for $\zeta\in\mu_2$.
For $n=2$ the two conditions on functions coincide. Let
$C_c^\infty(\overline{T(k_v)})_-$ and $C_c^\infty(\overline{T(\A)^1})_-$
be the spaces of anti-genuine functions that are smooth and compactly
supported on these locally compact groups: locally constant in the
non-archimedean variables and $C^\infty$ in the archimedean ones. As in
Section~\ref{sec:functions}, this is a condition on the topological
group, not on the function $t\mapsto f(t,1)$. For the split torus the two
agree, because $a\mapsto(a,1)$ is continuous on $A(k_v)$ and the
cocycle~\eqref{eq:cocycle-A} is locally constant
(Lemma~\ref{lem:local}(i)); for an anisotropic torus they do not, for
example for $SO(2)\subset SL(2,\R)$, whose preimage is connected
(Lemma~\ref{lem:local}(d)). At places where $T(\OO_v)^*$ is defined
let $e_{T,v}$ be the anti-genuine function equal to $1$ on
$T(\OO_v)^*$ and supported on $\overline{T(\OO_v)}$. Let
$\mathcal H(\overline{T(\A)})$ be the span of the products
$\bigotimes_vf_v$, $f_v\in C_c^\infty(\overline{T(k_v)})_-$, with
$f_v=e_{T,v}$ for almost all $v$; these descend to $\overline{T(\A)}$ as
in Section~\ref{sec:global}. Measures on $T(\A)$ are as in
Section~\ref{sec:global}, and on covers $\mu_2$ has total volume $1$.

For a locally compact group $H$ let $\Pi(H)$ be the set of
(equivalence classes of) irreducible admissible genuine representations
of $H$, and $\Pi_u(H)$ the unitary ones.

\begin{lemma}\label{lem:genuine-chars}
Let $F=k_v$.
\begin{enumerate}[label=(\alph*)]
\item Every element of $\Pi(\overline{T(F)})$, $\Pi(\overline{T(\A)})$ or
  $\Pi(\overline{T(\A)^1})$ is a genuine character.
\item Genuine characters of $\overline{T(F)}$ exist, and any two differ
  by a character of $T(F)$: they form a torsor under $T(F)^\wedge$.
\item Every genuine character $\bar\chi$ of $\overline{T(\A)}$ factors
  uniquely as $\bar\chi=\prod_v\bar\chi_v$, where $\bar\chi_v$ is a
  genuine character of $\overline{T(k_v)}$, trivial on $T(\OO_v)^*$ for
  almost all $v$. Conversely, every such family defines a genuine
  character of $\overline{T(\A)}$. The same holds for unitary characters.
\end{enumerate}
\end{lemma}

\begin{proof}
(a) The groups are abelian, so Schur's lemma applies; for unitary
representations of locally compact abelian groups this is standard, and
for admissible representations Schur's lemma holds because they have
countable dimension. Since $\mu_2$ is central, it acts by a character,
and that character is the identity because the representation is
genuine.

(b) $\mu_2$ is a closed subgroup of the locally compact abelian group
$\overline{T(F)}$, so the character $\zeta\mapsto\zeta$ of $\mu_2$
extends~\cite{Ro}. The quotient of two genuine characters is trivial on
$\mu_2$, hence is the pull-back of a character of $T(F)$.

(c) Put $\bar\chi_v=\bar\chi\circ\iota_v$, where $\iota_v$ is the
inclusion of $\overline{T(k_v)}$. Since $\bar\chi$ is continuous, there
is a neighborhood $W$ of the identity in $\overline{T(\A)}$ with
$|\bar\chi(x)-1|<1$ for $x\in W$. By the definition of the topology, $W$
contains the image of $\prod_{v\notin S}T(\OO_v)^*$ for some finite set
$S$ of places. That image is a subgroup, so its image under $\bar\chi$ is
a subgroup of $\C^\times$ contained in the disc $\{|z-1|<1\}$, and the
only such subgroup is $\{1\}$. Hence $\bar\chi_v$ is trivial on
$T(\OO_v)^*$ for $v\notin S$. (The kernel of $\bar\chi$ need not be open:
at an archimedean place $\bar\chi_v$ can be nontrivial on every
neighborhood of the identity.) Every $x\in\overline{T(\A)}$ is the image of
some $(x_v)$ with $x_v\in T(\OO_v)^*$ for $v\notin S'$, and then
$\bar\chi(x)=\prod_{v\in S'}\bar\chi_v(x_v)$. Conversely, $\prod_v
\bar\chi_v$ is a character of the restricted product, and it is trivial
on $\{(\zeta_v)\mid\prod\zeta_v=1\}$ because each $\bar\chi_v$ is
genuine. So it descends to $\overline{T(\A)}$.
\end{proof}

Let $X_T=T(\Q)\backslash\overline{T(\A)^1}$. It is a compact abelian
group, an extension of the compact group $T(\Q)\backslash T(\A)^1$ by
$\mu_2$; $\mu_2$ injects into $X_T$ because $T(\Q)$ meets $\mu_2$
trivially in $\overline{G(\A)}$. Let $\Pi_{\mathrm{aut}}(X_T)$ be the set
of genuine characters of $\overline{T(\A)^1}$ that are trivial on
$T(\Q)$, i.e.\ the genuine characters of $X_T$.

\begin{theorem}\label{thm:corr}
\begin{enumerate}[label=(\alph*)]
\item $L^2(X_T)_{\mathrm{gen}}=\widehat\bigoplus_{\bar\chi\in
  \Pi_{\mathrm{aut}}(X_T)}\C\,\bar\chi$, each character occurring with
  multiplicity one.
\item $\Pi_{\mathrm{aut}}(X_T)$ is not empty. For any
  $\bar\chi_0\in\Pi_{\mathrm{aut}}(X_T)$, the map
  $\chi\mapsto\bar\chi_0\cdot(\chi\circ p)$ is a bijection from the
  characters of $T(\Q)\backslash T(\A)^1$ onto $\Pi_{\mathrm{aut}}(X_T)$.
  Replacing $\bar\chi_0$ by another element changes this bijection by a
  translation of $(T(\Q)\backslash T(\A)^1)^\wedge$.
\item For $T=A$, the character $\chi_\psi$ of Lemma~\ref{lem:weil} lies
  in $\Pi_{\mathrm{aut}}(X_A)$; it depends only on the fixed additive
  character $\psi$ of $\A_k/k$.
\end{enumerate}
\end{theorem}

\begin{proof}
(a) Peter--Weyl for the compact abelian group $X_T$: $L^2(X_T)$ is the
Hilbert sum of its characters, each once, and the genuine subspace is
spanned by the genuine ones. (b) Apply Lemma~\ref{lem:genuine-chars}(b)
to the closed subgroup $\mu_2$ of $X_T$; the second sentence follows
as in that lemma. (c) is Lemma~\ref{lem:weil}(c).
\end{proof}

The bijection in~(b) depends on the choice of $\bar\chi_0$, and for a
general torus there is no canonical one. The choice cannot be avoided:

\begin{example}\label{ex:order4}
Let $k=\Q$, and let $\varepsilon\in A(\A)^1$ be $a(-1)$ at the real place
and $1$ at every finite place. By Lemma~\ref{lem:local}(e) any lift
$\tilde\varepsilon$ satisfies $\tilde\varepsilon^2=(1,(-1,-1)_\infty)
=(1,-1)$. So every genuine character $\bar\chi$ of $\overline{A(\A)^1}$
has $\bar\chi(\tilde\varepsilon)=\pm i$, and no genuine character is of
the form $(\chi\circ p)\cdot\mathrm{sgn}$ for a splitting $A(\A)^1\to
\overline{A(\A)^1}$; indeed no such splitting exists.
\end{example}

For the split torus there is a canonical choice, given by Weil
indices~\cite{We64}. Let $\psi=\prod_v\psi_v$ be the nontrivial
character of $\A_k/k$ fixed in Section~\ref{sec:global}. For $a\in
k_v^\times$ let $\gamma_v(a)\in\mu_8$ be the Weil index of the character
$x\mapsto\psi_v(ax^2)$ of $k_v$, normalized by
\begin{equation}\label{eq:weil-index}
\int_{k_v}\psi_v(ax^2)\,dx=\gamma_v(a)\,|2a|_v^{-1/2}
\end{equation}
(self-dual measure; for $v$ finite the integral is the limit over the
lattices $\varpi^{-M}\OO_v$, $M\to\infty$).

\begin{lemma}[Weil~\cite{We64}]\label{lem:weil}
\begin{enumerate}[label=(\alph*)]
\item $\gamma_v(a)$ depends only on the class of $a$ in
  $k_v^\times/k_v^{\times2}$, and
  $\gamma_v(ab)\gamma_v(1)=\gamma_v(a)\gamma_v(b)(a,b)_v$.
\item If $v$ is finite and non-dyadic and $\psi_v$ is unramified, then
  $\gamma_v(u)=\gamma_v(1)=1$ for $u\in\OO_v^\times$. For $a\in k^\times$,
  $\prod_v\gamma_v(a)=1$.
\item Hence $\chi_{\psi,v}(a(t),\zeta)=\zeta\,\gamma_v(t)/\gamma_v(1)$ is
  a genuine character of $\overline{A(k_v)}$, trivial on $A(\OO_v)^*=
  A(\OO_v)\times\{1\}$ for almost all $v$, and
  $\chi_\psi=\prod_v\chi_{\psi,v}$ is a genuine character of
  $\overline{A(\A)}$ trivial on $A(\Q)$. Its restriction to
  $\overline{A(\A)^1}$ lies in $\Pi_{\mathrm{aut}}(X_A)$.
\item For $b\in k^\times$ and $\psi_b(x)=\psi(bx)$,
  $\chi_{\psi_b}=\chi_\psi\cdot\bigl(a(t)\mapsto\prod_v(b,t_v)_v\bigr)$.
\end{enumerate}
\end{lemma}

\begin{proof}
(a) and (b) are due to Weil~\cite[\S\S14, 28]{We64}. (c) By~(a),
$\chi_{\psi,v}(a(t_1),1)\chi_{\psi,v}(a(t_2),1)=(t_1,t_2)_v\,
\chi_{\psi,v}(a(t_1t_2),1)=\chi_{\psi,v}((a(t_1),1)(a(t_2),1))$, using the
group law~\eqref{eq:cocycle-A}. Unramifiedness is~(b), and triviality on
$(a(t),1)$, $t\in k^\times$, is the product formula in~(b) applied to
$\gamma(t)$ and to $\gamma(1)$. (d) $\psi_{b,v}(ax^2)=\psi_v(abx^2)$, so
$\gamma_{\psi_b,v}(t)/\gamma_{\psi_b,v}(1)=\gamma_v(bt)/\gamma_v(b)=
(\gamma_v(t)/\gamma_v(1))\,(b,t)_v$ by~(a).
\end{proof}

\begin{example}\label{ex:weil-Q}
Let $k=\Q$, $\psi_p(x)=e^{2\pi i\{x\}_p}$ and $\psi_\infty(x)=
e^{-2\pi ix}$, so $\psi$ is trivial on $\Q$. Then
$\gamma_\infty(a)=e^{-\pi i\,\mathrm{sgn}(a)/4}$, and for odd $p$,
$\gamma_p(u)=1$ and $\gamma_p(pu)=\bigl(\tfrac up\bigr)\epsilon_p$ for
$u\in\Z_p^\times$, where $\epsilon_p=1$ if $p\equiv1\pmod4$ and
$\epsilon_p=i$ if $p\equiv3\pmod4$. So
$\chi_{\psi,\infty}(a(-1),1)=i$, in accordance with
Example~\ref{ex:order4}, and $\chi_{\psi,p}(a(p),1)=\epsilon_p$ for odd
$p$. The dyadic values are tabulated by the check script, which computes
$\gamma_v$ from~\eqref{eq:weil-index} as a finite Gauss sum and verifies
Lemma~\ref{lem:weil}(a)--(c) at $v=\infty,2,3,5,7$ and the product
formula on random $a\in\Q^\times$.
\end{example}

\subsection{The trace formula}\label{sec:tf-rank0}

Let $R$ be the right regular representation of $\overline{T(\A)^1}$ on
$L^2(X_T)_{\mathrm{gen}}$, and for $f\in C_c^\infty(\overline{T(\A)^1})_-$
put $R(f)=\int_{\overline{T(\A)^1}}f(y)R(y)\,dy$ and
\[
\bar\chi(f)=\int_{\overline{T(\A)^1}}f(y)\bar\chi(y)\,dy,\qquad
\bar\chi\in\Pi(\overline{T(\A)^1}).
\]
The integrand $f\bar\chi$ is invariant under $\mu_2$, so $\bar\chi(f)$ is
an integral over $T(\A)^1$; if $\bar\chi$ were trivial on $\mu_2$ instead,
the integral would vanish.

\begin{theorem}\label{thm:tf-rank0}
Let $f\in C_c^\infty(\overline{T(\A)^1})_-$.
\begin{enumerate}[label=(\alph*)]
\item $R(f)$ is the integral operator on $L^2(X_T)_{\mathrm{gen}}$ with
  kernel
  \[
  K_f(x,y)=\sum_{\delta\in T(\Q)}f(x^{-1}\delta y),
  \]
  and $R(f)\bar\chi=\bar\chi(f)\,\bar\chi$ for
  $\bar\chi\in\Pi_{\mathrm{aut}}(X_T)$.
\item $R(f)$ is of trace class, and
  \begin{equation}\label{eq:tf-rank0}
  \sum_{\bar\chi\in\Pi_{\mathrm{aut}}(X_T)}\bar\chi(f)
  =\operatorname{tr}R(f)=\int_{X_T}K_f(x,x)\,dx
  =\mathrm{vol}\bigl(T(\Q)\backslash T(\A)^1\bigr)\sum_{\delta\in T(\Q)}f(\delta),
  \end{equation}
  the sum on the left converging absolutely.
\end{enumerate}
\end{theorem}

\begin{proof}
(a) For $\varphi\in L^2(X_T)_{\mathrm{gen}}$,
\[
(R(f)\varphi)(x)=\int_{\overline{T(\A)^1}}f(y)\varphi(xy)\,dy
=\int_{\overline{T(\A)^1}}f(x^{-1}y)\varphi(y)\,dy
=\int_{X_T}\sum_{\delta\in T(\Q)}f(x^{-1}\delta y)\varphi(y)\,dy,
\]
because $\varphi(\delta y)=\varphi(y)$. For $\varphi=\bar\chi$ the first
expression is $\bar\chi(x)\bar\chi(f)$.

(b) The group is abelian, so $f(x^{-1}\delta x)=f(\delta)$, and
$\int_{X_T}K_f(x,x)\,dx=\mathrm{vol}(X_T)\sum_\delta f(\delta)$ with
$\mathrm{vol}(X_T)=\mathrm{vol}(T(\Q)\backslash T(\A)^1)$ since $\mu_2$
has volume $1$. The function $F(x)=\sum_{\delta}f(\delta x)$ is smooth on
$X_T$ (the sum is locally finite, $T(\Q)$ being discrete), so its Fourier
series converges absolutely; its Fourier coefficient at $\bar\chi$ is
$\int_{X_T}F\bar\chi=\bar\chi(f)$, which vanishes unless $\bar\chi$ is
genuine. Evaluating the Fourier series at $x=1$ gives the outer equality
of~\eqref{eq:tf-rank0}, i.e.\ Poisson summation. Since $R(f)$ is diagonal
in the orthonormal basis $\mathrm{vol}(X_T)^{-1/2}\bar\chi$ with summable
eigenvalues $\bar\chi(f)$, it is of trace class with that trace.
\end{proof}

Every conjugacy class in $\overline{T(\A)^1}$ is a point and every
centralizer is the whole group, so the geometric side has no orbital
integrals, only the values $f(\delta)$. Here $f(\delta)$ means $f$ at
the image $(\delta,\kappa_\A(\delta))$ of $\delta$; for $T=A$ it is
$f(a(\delta),1)$.

The trace formula for $\overline{T(\A)^1}$ is literally Poisson summation
on $T(\Q)\backslash T(\A)^1$, after multiplying by a genuine automorphic
character:

\begin{lemma}\label{lem:transfer}
Fix $\bar\chi_0\in\Pi_{\mathrm{aut}}(X_T)$ with local components
$\bar\chi_{0,v}$, as in Lemma~\ref{lem:genuine-chars}(c). For an
anti-genuine $\bar f$ on $\overline{T(\A)^1}$ (resp.\
$\overline{T(k_v)}$) define $f$ on $T(\A)^1$ (resp.\ $T_k(k_v)$) by
$f(t)=\bar f(\tilde t)\bar\chi_0(\tilde t)$ (resp.\ $\bar\chi_{0,v}$),
$\tilde t$ any lift of $t$.
\begin{enumerate}[label=(\alph*)]
\item $\bar f\mapsto f$ is well defined, bijective from
  $C_c^\infty(\overline{T(\cdot)})_-$ onto $C_c^\infty(T(\cdot))$, and
  carries convolution to convolution.
\item At every $v$ where $\bar\chi_{0,v}$ is trivial on $T(\OO_v)^*$
  (almost all $v$), it carries $e_{T,v}$ to the characteristic function of
  $T_k(k_v)\cap K_v$, which is the unit of the Hecke algebra of
  $T_k(k_v)$ when $T_k(k_v)\cap K_v$ has volume $1$.
\item $f(\delta)=\bar f(\delta)$ for $\delta\in T(\Q)$, and
  $(\bar\chi_0\cdot(\chi\circ p))(\bar f)=\hat f(\chi):=\int_{T(\A)^1}
  f(t)\chi(t)\,dt$ for every character $\chi$ of $T(\A)^1$.
\end{enumerate}
Consequently~\eqref{eq:tf-rank0} for $\bar f$ is the Poisson summation
formula for $f$,
\[
\sum_\chi\hat f(\chi)=\mathrm{vol}\bigl(T(\Q)\backslash T(\A)^1\bigr)
\sum_{\delta\in T(\Q)}f(\delta),
\]
with $\chi$ running over the characters of $T(\Q)\backslash T(\A)^1$.
\end{lemma}

\begin{proof}
(a) $\bar f\bar\chi_0$ is $\mu_2$-invariant because $\bar f$ is
anti-genuine and $\bar\chi_0$ genuine, and it is smooth because
$\bar\chi_0$ is a continuous character; so it is the pull-back of a
smooth function on $T$. The inverse map is
$f\mapsto(f\circ p)\bar\chi_0^{-1}$. For convolution,
$(\bar f_1*\bar f_2)(x)\bar\chi_0(x)=\int\bar f_1(y)\bar\chi_0(y)\,
\bar f_2(y^{-1}x)\bar\chi_0(y^{-1}x)\,dy$, and the integrand is
$\mu_2$-invariant in $y$, so the integral over the cover equals the
integral over $T$. (b) is immediate. (c) $\bar\chi_0$ is trivial on
$T(\Q)$; the second identity is the definition, the integrand being
$\mu_2$-invariant.
\end{proof}

Together with Theorem~\ref{thm:tf-rank0}, the lemma shows that the
correspondence $\chi\leftrightarrow\bar\chi_0\cdot(\chi\circ p)$ of
Theorem~\ref{thm:corr} matches the two trace formulas term by term.

\subsection{Distributions for Section 6}\label{sec:IA}

Let $S$ be a finite set of places of $\Q$ and let $\overline{A(\A_S)}$ be
the cover of $A(\A_S)=\prod_{p\in S}A(\Q_p)$ defined by
$\prod_{p\in S}\beta_p$. Let $\mathcal H(\overline{A(\A_S)})$ be the
algebra of functions $\bigotimes_{p\in S}f_p$, $f_p\in
C_c^\infty(\overline{A(\Q_p)})_-$, and the analogous algebra
$\mathcal H(\overline{A(\A)})$ of Section~\ref{sec:torus-cover}. For a
class $\mathcal O=\{\delta\}$ of $A(\Q)$ (a single element, since the
cover is abelian) and $\bar\chi\in\Pi_u(\overline{A(\A_S)})$, put
\begin{equation}\label{eq:IA}
I^A_{\mathcal O}(f)=f(a(\delta),1),\qquad
I^A_{\bar\chi}(f)=\bar\chi(f)=\int_{\overline{A(\A_S)}}f(y)\bar\chi(y)\,dy,
\end{equation}
for $f\in\mathcal H(\overline{A(\A_S)})$, and likewise over $\A$. Here
$(a(\delta),1)$ is the image of $\delta$ under the rational splitting
(Lemma~\ref{lem:global}(b)). Both are invariant distributions: the cover
is abelian, so conjugation acts trivially.

\begin{remarkrev}\label{rem:general-n-rank0}
For general $n$ the cover of $A$ is not abelian: by the computation in
Lemma~\ref{lem:local}(f), $\overline{A(F)}$ is a Heisenberg-type group
whose commutator pairing is $(t_1,t_2)_n^{-2}$. Its genuine irreducible
representations are then parametrized by genuine characters of its
center, the preimage of $A^N$ with $N=n/\gcd(n,2)$. Matching them with
characters of $A^N$ requires a splitting over $A^N$, and the global
argument requires a Hasse principle for $N$-th powers, which holds when
$\mu_N\subset k$. These are referee items II.1--II.4; none of them arises
for $n=2$, and this section does not treat $n>2$.
\end{remarkrev}

\subsection{Status of the results}\label{sec:status3}

\begin{remarkrev}\label{rem:sec3-status}
\emph{What rests on citation.} Taken from the literature:
Lemma~\ref{lem:weil}(a)--(b)~\cite{We64}; the compactness of
$T(\Q)\backslash T(\A)^1$~\cite{PR}; and the extension of characters from
closed subgroups of locally compact abelian groups~\cite{Ro}. Everything
else in this section is proved in the text.
\end{remarkrev}

\begin{remarkrev}\label{rem:verified-rank0}
The script \texttt{checks/sec3\_n2\_check.py} verifies: that the 1995
cocycle on $A$ equals~\eqref{eq:cocycle-A}; that $\gamma_v$, computed
from~\eqref{eq:weil-index}, is an eighth root of unity, depends only on
the square class, is $1$ on units at $p=3,5,7$, and satisfies
Lemma~\ref{lem:weil}(a); that $\chi_{\psi,v}$ respects the group law
(Lemma~\ref{lem:weil}(c)) at $v=\infty,2,3,5,7$; the product formula of
Lemma~\ref{lem:weil}(b) for random $a\in\Q^\times$; the odd-$p$ values of
Example~\ref{ex:weil-Q}; and Example~\ref{ex:order4}. The rest of the
section is proved in the text.
\end{remarkrev}
